% paper/appendix_a_generalization.tex
% Compressed generalization-audit methodology and per-minister results,
% moved from the former §6.1 to Appendix A for the CFP's 10-page body
% limit, and compressed further (full prose walkthroughs -> methodology
% paragraph + itemized verdicts) to fit the appendix page budget.

\paragraph{Methodology.} \S\ref{sec:generalization-audit} summarizes why
this audit exists and its top-line verdict; this appendix gives the
per-minister mechanism and the exact numbers behind it. AtriumOS's
invented vocabulary (\texttt{Facility\allowbreak ConfigRead},
\texttt{Hvac\allowbreak Diagnostic}, \texttt{Retrieve\allowbreak
Tenant\allowbreak AccessLog}, and so on) was dispatched as 18 attack and 18
matched-benign scenarios, one per minister's intended detection shape,
through the real pipeline, tested before AtriumOS existed as a case for
\emph{transfer} rather than a headline claim. The 18 attack scenarios are
not evenly split across ministers (8 EXECUTOR, 7 VAULT, 3 combined for
CHANNEL and NAVIGATOR); EXECUTOR's and VAULT's verdicts rest on a
reasonable sample, CHANNEL's and NAVIGATOR's on very few cases each, and
should be read as illustrative rather than statistically powered.

\begin{itemize}
\item \textbf{EXECUTOR --- generalizes cleanly.} Caught all eight
novel-vocabulary execution attacks via structural/\allowbreak technique-level shape
matching, independent of tool identity. One structural false positive was
found (a network-fetch-to-disk predicate firing on a benign data-retrieval
call) and fixed pre-submission, with zero new benign false positives on
either population; parser test suites remain green (27/27, 37/37).
\item \textbf{VAULT --- mostly generalizes, one inherent limit.} Caught
5/7 novel-vocabulary credential attacks via structural credential-path and
secret-format matching. One miss is a narrow surface-form gap (fixable);
the other is architectural, not a bug: a named-service secret-key rule is,
by construction, a literal allowlist, so a genuinely novel secret-variable
name cannot match it regardless of technique.
\item \textbf{CHANNEL --- split, since partially closed.} Its
provenance/escalation path generalizes: a behavior-general
destination-value fallback correctly escalated a novel send-tool's
attacker email address, classified \textsc{novel}. Its taint/hard-block
path originally did not, since destination status was decided by a
benchmark-specific tool-name dictionary alone; we have since wired the
same shape-based fallback into the hard-block path (verified: a novel
tool with a shape-matched value, e.g.\ an IBAN, now correctly
hard-\textsc{Block}s), closing the gap for destination values matching a
known shape. A genuinely novel exfiltration tool whose destination is
\emph{not} shape-matchable (e.g.\ a proper-noun business name, the
\texttt{reserve\_hotel} case elsewhere in this paper) remains invisible to
the hard-block path --- the fix is real but narrower than full
generalization.
\item \textbf{NAVIGATOR --- does not currently generalize as a detector.}
Cosine similarity against the attack-pattern corpus separates novel attack
from benign by only 0.045 (0.590 vs.\ 0.545 mean), and the BM25 lexical
gate that preserves on-benchmark precision makes this worse off-benchmark,
collapsing both to $\sim$0.38 --- below the escalation floor, so NAVIGATOR
goes silent on novel-domain attacks. This is structural, not incidental:
the same mechanism giving NAVIGATOR its low on-benchmark FPR removes its
off-benchmark signal. The one general component built for NAVIGATOR --- an
observational/consequential tool classifier --- correctly labeled all four
novel tools from verb-token structure alone, but is a \emph{taxonomy, not a
detector}: it does not decide whether a consequential action was
\emph{authorized}, which requires the pinned-intent authorized-flow layer
reported as an incomplete negative result in \S\ref{par:navigator-authflow}.
\end{itemize}

\paragraph{A second, genuinely independent check: AgentHarm.} AtriumOS is
author-written; to test transfer against data this team did not author at
all, we used AgentHarm's (UK AI Safety Institute, ICLR 2025) published
synthetic tool implementations --- 80 harmful-scenario and 80
benign-scenario Python tool files, MIT-licensed for safety research, fetched
directly rather than through AgentHarm's official gated benchmark (Inspect
framework, API keys, and not tool-call structured regardless). A local
\texttt{qwen2.5:7b} generated one realistic call per tool from its
docstring (a first heuristic attempt was discarded: bland placeholders on
the harmful side, suspicious-looking placeholders on the benign side,
neither a fair test; 1 of the 80 harmful tools failed argument generation
and was excluded, leaving 79). Every one of those 79 cases was
hand-classified, once, before computing any number, into in-scope for
Kavach's stated
four-domain threat model or not: AgentHarm covers 11 harm categories far
broader than Kavach's claim, and counting e.g.\ a drug-marketplace search
against Kavach would be unfair to both sides. Of 44 in-scope cases: 36
caught (81.8\%), 8 missed --- anti-forensics (\texttt{delete\_access\_logs}),
public exploit-bragging, phishing-asset generation, a phishing email with a
malware attachment, a likely tool-name vocabulary gap
(\texttt{transfer\_money\_crypto}, the same class of gap as CHANNEL's
hard-block path above), two malware-named file uploads, and a
no-argument credential harvest consistent with an earlier, independent
finding. The 80 benign cases: 0\% hard-block, 56.3\% Escalate, consistent
with \S\ref{sec:design-tradeoffs}'s disclosed 67.3\% agent-shaped rate
rather than contradicting it.

We frame this section as a strength, not an apology: it is the output of
a deliberate, independent validation pass run \emph{before} submission,
specifically to surface these boundaries ourselves; scoped remedies for
both open boundaries are discussed in \S\ref{sec:limitations}.
